\documentclass[11pt]{article}

\usepackage[margin=1.0in]{geometry}
\usepackage{amsmath,amssymb,amsthm,mathtools,array,booktabs}
\usepackage[hidelinks]{hyperref}
\usepackage{microtype}

\numberwithin{equation}{section}

\newcommand{\A}{\alpha'}
\newcommand{\dd}{\mathrm d}
\newcommand{\ee}{\mathrm e}
\newcommand{\ii}{\mathrm i}
\newcommand{\cA}{\mathcal A}
\newcommand{\cB}{\mathcal B}
\newcommand{\cD}{\mathcal D}
\newcommand{\cF}{\mathcal F}
\newcommand{\cH}{\mathcal H}
\newcommand{\cI}{\mathcal I}
\newcommand{\cK}{\mathcal K}
\newcommand{\cM}{\mathcal M}
\newcommand{\cN}{\mathcal N}
\newcommand{\cV}{\mathcal V}
\newcommand{\cZ}{\mathcal Z}
\newcommand{\Tr}{\operatorname{Tr}}
\newcommand{\im}{\operatorname{Im}}
\newcommand{\vol}{\operatorname{Vol}}
\newcommand{\wh}{\widehat}
\newcommand{\wt}{\widetilde}
\newcommand{\abs}[1]{\left|#1\right|}
\newcommand{\vev}[1]{\left\langle #1\right\rangle}

\newtheorem{minimalinput}{Minimal input}
\newtheorem{warning}{Convention warning}

\title{Absolute Normalization of the Compact $c=1$ String\\
Ingredient by Ingredient in Xi's Worldsheet Conventions}
\author{Scratch derivation from the action, CFT data, and sewing unitarity}
\date{21 July 2026}

\begin{document}
\maketitle

\begin{abstract}
We derive the genus-one and genus-two vacuum integrands of the compact $c=1$
bosonic string using the logic of the string notes.  Every ingredient is
computed separately: the compact zero mode and lattice action, oscillator
determinants, Liouville spectral measure, ghost determinant, moduli form,
loop-momentum Gaussian, Mumford-form residue, topology sewing constant, and
global stack weight.  The matrix model is not used as a normalization input.
It is introduced only after the worldsheet answer has been fixed, as an
external comparison.

The main bookkeeping distinction is between the universal BRST normalization
$K_{S^2}^{\rm crit}=8\pi/\A$ used in the critical-string measure and the
reduced physical $c=1$ sphere normalization
$\wh K_{S^2}^{c=1}=2\pi\sqrt{\A}\,\wt K_{S^2}=4\pi$.  Sewing then gives
\[
 \frac{C_g^{c=1}}{C_g^{\rm crit}}
 =\left(\frac{2}{\A}\right)^{g-1}.
\]
The torus vacuum is unchanged, while the genus-two local kernel is multiplied
by $2/\A$.  At $\A=1$ the positive-real, $g_s^2$-stripped coefficient is
$2/\pi$.
\end{abstract}

\tableofcontents

\section{Logic contract: what is input and what is derived}

The purpose of this note is to prevent an overall constant from entering
through an implicit convention.  Every numerical factor must belong to one of
the following categories.

\begin{minimalinput}[Worldsheet theory]
The matter theory is a timelike free boson $X^0$ times $b=1$ Liouville theory,
together with the standard $(b,c)$ ghosts and a constant dilaton $\Phi_0$.
We set $g_s=\ee^{\Phi_0}$.
\end{minimalinput}

\begin{minimalinput}[Ghost normalization]
On the sphere,
\begin{equation}
 \vev{\prod_{i=1}^3c(z_i)\widetilde c(\bar z_i)}_{S^2}
 =\abs{z_{12}z_{13}z_{23}}^2.
 \label{eq:ghost-sphere}
\end{equation}
This is Xi's equation (4.32).
\end{minimalinput}

\begin{minimalinput}[Liouville states]
The reflection-symmetric Liouville primary $V_P$, $P\geq0$, is normalized by
\begin{equation}
 \vev{V_P(1)V_{P'}(0)}=\pi\delta(P-P'),
 \qquad
 \mathbf 1_{\cH_L}=\int_0^\infty\frac{\dd P}{\pi}\,
 |P\rangle\langle P|,
 \label{eq:liouville-metric}
\end{equation}
and
\begin{equation}
 \vev{V_{P_1}V_{P_2}V_{P_3}}=C(P_1,P_2,P_3).
 \label{eq:liouville-three}
\end{equation}
These are Xi's equations (I.39)--(I.42).
\end{minimalinput}

\begin{minimalinput}[Physical string states]
The one-string in/out states and their vertices are
\begin{align}
 {}_{\rm in/out}\!\langle\omega|\omega'\rangle_{\rm in/out}
 &=\omega\delta(\omega-\omega'),
 \label{eq:state-metric}\\
 \cV_\omega^\pm
 &=g_s\,c\widetilde c\,
 \exp\!\left(\pm\frac{\ii\omega X^0}{\sqrt{\A}}\right)
 V_{P=\omega/2}.
 \label{eq:physical-vertex}
\end{align}
These are Xi's equations (4.111)--(4.112).
\end{minimalinput}

\begin{minimalinput}[Plumbing coordinate and differential forms]
Local coordinates are sewn literally by
\begin{equation}
 uv=q,\qquad q=\ee^{2\pi\ii\tau}.
 \label{eq:plumbing}
\end{equation}
The cylinder has circumference $2\pi$.  Differential forms are ordered as in
the string notes; in particular
\begin{equation}
 \ii\,\dd z\wedge\dd\bar z=2\,\dd(\Re z)\dd(\Im z).
 \label{eq:complex-real}
\end{equation}
\end{minimalinput}

\begin{minimalinput}[Unitarity]
With $S=1+\ii T$, impose $SS^\dagger=1$.  A one-particle pole must factorize
with the inverse of the physical two-point metric.  On the worldsheet this pole
is produced by the region $q\to0$ of moduli space.
\end{minimalinput}

No matrix-model free energy, no numerical Monte Carlo result, and no guessed
integer or power of $2$ will be used below.

\section{Worldsheet conventions}

\subsection{Action and central charge}

In Lorentzian target signature, the matter action is
\begin{align}
 S_{X^0}
 &=-\frac{1}{4\pi\A}\int_\Sigma\dd^2\sigma\sqrt g\,
 g^{ab}\partial_aX^0\partial_bX^0,
 \label{eq:timelike-action}\\
 S_L[X^1]
 &=\frac{1}{4\pi\A}\int_\Sigma\dd^2\sigma\sqrt g\,
 \left[
 g^{ab}\partial_aX^1\partial_bX^1
 +\sqrt{\A}\,QX^1R(g)
 +4\pi\mu_L\ee^{2bX^1/\sqrt{\A}}
 \right].
 \label{eq:liouville-action}
\end{align}
Here $Q=b+b^{-1}$.  At $b=1$, $Q=2$ and
\begin{equation}
 c_{X^0}+c_L+c_{bc}=1+25-26=0.
\end{equation}
The constant dilaton contributes
\begin{equation}
 S_{\Phi_0}=\Phi_0\chi(\Sigma),\qquad
 \ee^{-S_{\Phi_0}}=g_s^{2g-2}.
 \label{eq:dilaton-weight}
\end{equation}
Since each physical vertex in \eqref{eq:physical-vertex} contains one $g_s$,
a genus-$g$, $n$-point amplitude is of order $g_s^{2g-2+n}$.

For the thermal vacuum amplitude we Wick rotate $X^0$ to a Euclidean compact
boson $X$ with
\begin{equation}
 X\sim X+2\pi R_{\rm phys},
 \qquad r:=\frac{R_{\rm phys}}{\sqrt{\A}}.
 \label{eq:target-circle}
\end{equation}
The worldsheet torus coordinate has periods $2\pi$ and $2\pi\tau$.

The action also fixes the conformal weights used in the sewing calculation.
With the exponential in \eqref{eq:physical-vertex},
\begin{equation}
 h_{X^0}=\bar h_{X^0}=-\frac{\omega^2}{4},
 \qquad
 h_L=\bar h_L=\frac{Q^2}{4}+P^2=1+P^2
 \quad (b=1).
 \label{eq:weights}
\end{equation}
Thus the product is a $(1,1)$ primary precisely when $P=\omega/2$.  Off this
physical locus the quantity that appears in a plumbing propagator is
\begin{equation}
 a(P,\omega)=P^2-\frac{\omega^2}{4}.
 \label{eq:offshell-a}
\end{equation}

\subsection{The special \texorpdfstring{$b=1$}{b=1} three-point identity}

The DOZZ coefficient in the normalization \eqref{eq:liouville-metric} obeys
\begin{equation}
 C(P_1,P_2,P_1+P_2)=8P_1P_2(P_1+P_2).
 \label{eq:resonance}
\end{equation}
On the physical kinematic locus $P_i=\omega_i/2$ and
$\omega=\omega_1+\omega_2$, this becomes
\begin{equation}
 C\!\left(\frac\omega2,\frac{\omega_1}{2},\frac{\omega_2}{2}\right)
 =\omega\omega_1\omega_2.
 \label{eq:energy-polynomial}
\end{equation}
This identity fixes the local three-string interaction once the sphere zero
mode and the string coupling are normalized.

\section{Unitarity and the sphere constants}

\subsection{Critical BRST state metric}

Let $N_{g,n}$ multiply Xi's top differential form on $\cM_{g,n}$.  Plumbing
two surfaces of types $(g_1,n_1)$ and $(g_2,n_2)$ and matching the pole to the
one-particle unitarity sum gives
\begin{equation}
 N_{g_1,n_1}N_{g_2,n_2}
 =-\ii N_{g,n}\frac{8\pi}{\A K_{S^2}^{\rm crit}},
 \qquad
 g=g_1+g_2,\qquad n=n_1+n_2-2.
 \label{eq:separating-recurrence}
\end{equation}
Pinching one handle gives
\begin{equation}
 N_{g-1,n+2}
 =-\ii N_{g,n}\frac{8\pi}{\A K_{S^2}^{\rm crit}}.
 \label{eq:nonseparating-recurrence}
\end{equation}
These are Xi's equations (4.69)--(4.70).  The simultaneous solution adopted in
the notes is
\begin{equation}
 \boxed{
 N_{g,n}=\ii^{3g-3+n},
 \qquad
 K_{S^2}^{\rm crit}=\frac{8\pi}{\A}.}
 \label{eq:critical-solution}
\end{equation}
To see that the numerical value of $K_{S^2}^{\rm crit}$ is fixed rather than
assumed, substitute the phase ansatz into the separating recurrence.  Since
$g=g_1+g_2$ and $n=n_1+n_2-2$,
\begin{equation}
 \frac{N_{g_1,n_1}N_{g_2,n_2}}{N_{g,n}}
 =\ii^{-1}=-\ii.
 \label{eq:N-ratio}
\end{equation}
Comparison with \eqref{eq:separating-recurrence} therefore requires
$8\pi/(\A K_{S^2}^{\rm crit})=1$, which is exactly the second result in
\eqref{eq:critical-solution}.  The nonseparating recurrence gives the same
condition.
In particular,
\begin{equation}
 N_{1,0}=1,\qquad N_{2,0}=-\ii.
 \label{eq:N10N20}
\end{equation}

\subsection{Physical \texorpdfstring{$c=1$}{c=1} timelike sphere normalization}

In the $c=1$ scattering convention, write the timelike zero-mode correlator as
\begin{equation}
 \vev{\ee^{-\ii k^0X^0}}_{S^2,X^0}
 =2\pi\ii\,\delta(k^0)\,g_s^{-2}\wt K_{S^2}.
 \label{eq:timelike-sphere}
\end{equation}
Here is the pole calculation behind the normalization.  The literal fixture
$uv=q$ gives, for the radial part of a closed-string propagator,
\begin{equation}
 \int_{|q|<1}\frac{\dd^2q}{2\pi|q|^2}|q|^{2a}
 =\frac{1}{2a},
 \qquad a=P^2-\frac{\omega^2}{4}.
 \label{eq:q-radial-pole}
\end{equation}
The Liouville completeness measure then has the on-shell discontinuity
\begin{equation}
 \operatorname{Im}\int_0^\infty\frac{\dd P}{\pi}
 \frac{1}{P^2-\omega^2/4-\ii0}
 =\int_0^\infty\dd P\,
 \delta\!\left(P^2-\frac{\omega^2}{4}\right)
 =\frac1\omega.
 \label{eq:liouville-cut}
\end{equation}
This is exactly the inverse metric required by
$\langle\omega|\omega'\rangle=\omega\delta(\omega-\omega')$.  In Xi's
closed-string propagator convention the remaining zero-mode residue is
$\sqrt{\A}\,\wt K_{S^2}/2$; the factor $1/2$ is the radial residue in
\eqref{eq:q-radial-pole}.  Requiring unit coefficient gives
\begin{equation}
 \frac{\sqrt{\A}\,\wt K_{S^2}}{2}=1,
\end{equation}
and hence
\begin{equation}
 \boxed{\wt K_{S^2}=\frac{2}{\sqrt{\A}}.}
 \label{eq:Ktilde}
\end{equation}
This is Xi's equation (4.114).  The momentum in
\eqref{eq:timelike-sphere} is physical,
\begin{equation}
 k^0=\frac{\omega}{\sqrt{\A}},
 \qquad
 \delta(k^0)=\sqrt{\A}\,\delta(\omega).
 \label{eq:delta-jacobian}
\end{equation}
Consequently the coefficient of the dimensionless energy delta function is
\begin{equation}
 \boxed{
 \wh K_{S^2}^{c=1}
 :=2\pi\sqrt{\A}\,\wt K_{S^2}
 =4\pi.}
 \label{eq:Khat}
\end{equation}

\begin{warning}
$K_{S^2}^{\rm crit}$ and $\wt K_{S^2}$ are not two names for the same
quantity.  The first normalizes the universal BRST state metric in the
critical-string derivation.  The second is the timelike zero-mode constant in
the physical $c=1$ scattering normalization.  After converting to the
dimensionless energy delta function, the $c=1$ topology metric is $\wh K$.
\end{warning}

\subsection{Tree-level \texorpdfstring{$1\to2$}{1 to 2} check}

Three vertices give $g_s^3$ and the sphere dilaton factor gives $g_s^{-2}$.
Using \eqref{eq:Khat} and \eqref{eq:energy-polynomial}, the coefficient of
$\delta(\omega-\omega_1-\omega_2)$ is
\begin{align}
 \cA_0(\omega_1,\omega_2)
 &=\ii g_s\,\wh K_{S^2}^{c=1}
 C\!\left(\frac\omega2,\frac{\omega_1}{2},\frac{\omega_2}{2}\right)
 \notag\\
 &=\boxed{4\pi\ii g_s\,\omega\omega_1\omega_2},
 \qquad \omega=\omega_1+\omega_2.
 \label{eq:tree-amplitude}
\end{align}
This reproduces Xi's equation (4.121) without using the matrix model.

\section{Topology constants from sewing}

First strip the single energy-conservation delta function and the local CFT
density from every connected amplitude.  Denote the remaining topology
constant by $C_g$.  With this reduced-amplitude convention, the tree result
\eqref{eq:tree-amplitude} says directly that
\begin{equation}
 C_0^{c=1}=\wh K_{S^2}^{c=1}=4\pi.
 \label{eq:C0-physical}
\end{equation}
Thus the Fourier $2\pi$ and the Jacobian in \eqref{eq:delta-jacobian} are part
of the reduced sphere coefficient, whereas the polynomial
$C(\omega/2,\omega_1/2,\omega_2/2)$ remains in the local CFT density.

Sewing a genus-$g_1$ and a genus-$g_2$ surface through one inverse two-point
metric gives
\begin{equation}
 C_{g_1+g_2}\,K_{S^2}=C_{g_1}C_{g_2}.
 \label{eq:topology-sewing}
\end{equation}
The canonical torus trace fixes $C_1=1$.  Taking $g_2=1$ recursively, or
solving \eqref{eq:topology-sewing} directly, gives
\begin{equation}
C_g=(K_{S^2})^{1-g},\qquad C_0=K_{S^2}.
 \label{eq:Cg}
\end{equation}
Therefore replacing the critical sphere metric by the physical $c=1$ metric
changes a genus-$g$ vacuum density by
\begin{equation}
 \boxed{
 \Lambda_g
 :=\frac{C_g^{c=1}}{C_g^{\rm crit}}
 =\left(\frac{K_{S^2}^{\rm crit}}{\wh K_{S^2}^{c=1}}\right)^{g-1}
 =\left(\frac{2}{\A}\right)^{g-1}.}
 \label{eq:Lambda-g}
\end{equation}
Two immediate consequences are
\begin{equation}
 \boxed{\Lambda_1=1,\qquad \Lambda_2=\frac{2}{\A}.}
 \label{eq:Lambda12}
\end{equation}

This is the precise sense in which a genus-one vacuum computation cannot see
the sphere-normalization mismatch: the torus has zero Euler characteristic.
A constant rescaling in the $g_s$--$\mu$ dictionary may shift a regulated
$\log\mu$ by an additive constant, but it cannot multiply the universal
coefficient of $\log\mu$.

\section{Genus one: derivation of every ingredient}

\subsection{Geometry and classical compact-boson sectors}

Let $\sigma^1,\sigma^2\in[0,2\pi)$ and
$z=\sigma^1+\tau\sigma^2$.  In these coordinates the flat torus metric is
\begin{equation}
 g_{ab}=\begin{pmatrix}1&\tau_1\\ \tau_1&|\tau|^2\end{pmatrix},
 \qquad
 \sqrt g=\tau_2,
 \qquad
 g^{ab}=\frac1{\tau_2^2}
 \begin{pmatrix}|\tau|^2&-\tau_1\\-\tau_1&1\end{pmatrix}.
 \label{eq:torus-metric}
\end{equation}
A map to the target circle can wind independently around the two cycles:
\begin{equation}
 X_{m,n}=x_0+R_{\rm phys}(m\sigma^1+n\sigma^2),
 \qquad m,n\in\mathbb Z.
\end{equation}
Substitution into the Euclidean version of \eqref{eq:timelike-action} gives
\begin{align}
 S_{m,n}
 &=\frac{(2\pi)^2}{4\pi\A}\,
 \tau_2 R_{\rm phys}^2(m,n)g^{-1}\binom mn \notag\\
 &=\frac{\pi R_{\rm phys}^2}{\A\tau_2}
 \left(|\tau|^2m^2-2\tau_1mn+n^2\right)
 =\frac{\pi R_{\rm phys}^2}{\A\tau_2}|m\tau-n|^2.
 \label{eq:compact-classical-action}
\end{align}
This calculation fixes both the $\pi$ and the placement of $\tau_2$ in the
lattice exponential.

\subsection{Compact zero mode and oscillator determinant}

The nonzero modes of one real scalar give
\begin{equation}
 Z_X^{\rm osc}(\tau)
 =\tau_2^{-1/2}|q|^{-1/12}
 \prod_{k=1}^\infty|1-q^k|^{-2}
 =\frac1{\sqrt{\tau_2}|\eta(\tau)|^2},
 \qquad q=\ee^{2\pi\ii\tau}.
 \label{eq:compact-oscillator}
\end{equation}
The zero-mode measure compatible with Xi's noncompact momentum measure
$\dd k/(2\pi)$ is
\begin{equation}
 \int_0^{2\pi R_{\rm phys}}\frac{\dd x_0}{2\pi\sqrt{\A}}
 =\frac{R_{\rm phys}}{\sqrt{\A}}.
 \label{eq:compact-zero-mode}
\end{equation}
This is not an additional convention: it is forced by the canonical trace.
Indeed, with
$p_{L,R}=n/R_{\rm phys}\pm wR_{\rm phys}/\A$, Poisson resummation uses
\begin{equation}
 \sum_{n\in\mathbb Z}
 \ee^{-\pi\tau_2\A n^2/R_{\rm phys}^2+2\pi\ii\tau_1nw}
 =\frac{R_{\rm phys}}{\sqrt{\A\tau_2}}
 \sum_{\widetilde n\in\mathbb Z}
 \ee^{-\pi R_{\rm phys}^2(\widetilde n-w\tau_1)^2/(\A\tau_2)},
 \label{eq:torus-poisson}
\end{equation}
which produces exactly the same zero-mode coefficient.  Combining
\eqref{eq:compact-classical-action}--\eqref{eq:compact-zero-mode},
\begin{equation}
 \boxed{
 Z_X(R_{\rm phys};\tau)
 =\frac{R_{\rm phys}}
 {\sqrt{\A\tau_2}|\eta(\tau)|^2}
 \sum_{m,n\in\mathbb Z}
 \exp\!\left[-\frac{\pi R_{\rm phys}^2}{\A\tau_2}
 |m\tau-n|^2\right].}
 \label{eq:compact-torus}
\end{equation}

\subsection{Liouville spectral trace}

The normalized Liouville identity resolution \eqref{eq:liouville-metric}
restricts the reflection-symmetric spectrum to $P\geq0$.  In Xi's dimensionful
momentum, $P_{\rm note}=2P/\sqrt\A$, the character is
$\ee^{-\pi\A\tau_2P_{\rm note}^2}/|\eta|^2$.  Therefore
\begin{align}
 \frac{Z_L(\tau)}{V_\phi}
 &=\int_0^\infty\frac{\dd P_{\rm note}}{\pi}
 \frac{\ee^{-\pi\A\tau_2P_{\rm note}^2}}
 {|\eta(\tau)|^2} \notag\\
 &=\frac1{\pi|\eta|^2}
 \frac{\sqrt\pi}{2\sqrt{\pi\A\tau_2}}
 =\boxed{\frac1{2\pi\sqrt{\A\tau_2}|\eta(\tau)|^2}}.
 \label{eq:liouville-torus}
\end{align}
The factor $1/2$ is thus the half-line Gaussian, not an automorphism factor.

\subsection{Ghosts, modulus, conformal Killing volume, and stack weight}

The chiral nonzero-mode trace of the weight-$(2,-1)$ ghost system is
$\eta(\tau)^2$; the left-right product is $|\eta|^4$.  The $b$ zero modes are
saturated by the Beltrami insertions for $\tau$ and $\bar\tau$.  With the
sphere ghost normalization \eqref{eq:ghost-sphere}, the resulting unmarked
torus formula is Xi's equation (4.90),
\begin{equation}
 \cA_1
 =\frac12\int_{\cF_1}
 \underbrace{\frac{\ii\,\dd\tau\wedge\dd\bar\tau}{2\tau_2}}
 _{\text{modulus and translation CKV}}
 \underbrace{|\eta(\tau)|^4}_{bc\ \text{nonzero modes}}
 Z_m(\tau,\bar\tau).
 \label{eq:torus-master}
\end{equation}
Here $N_{1,0}=1$ and $g_s^{2g-2}=1$.  The prefactor $1/2$ is the generic
$z\mapsto-z$ stabilizer of an unmarked torus when $\cF_1$ is a coarse
$PSL(2,\mathbb Z)$ fundamental domain.  It is distinct from the translation
CKV factor $1/(2\tau_2)$.  Finally,
\begin{equation}
 \ii\,\dd\tau\wedge\dd\bar\tau=2\,\dd^2\tau
 \quad\Longrightarrow\quad
 \frac{\ii\,\dd\tau\wedge\dd\bar\tau}{2\tau_2}
 =\frac{\dd^2\tau}{\tau_2}.
 \label{eq:torus-real-measure}
\end{equation}

\subsection{Multiplication: the complete genus-one integrand}

Insert $Z_m=Z_XZ_L$ into \eqref{eq:torus-master}.  Every factor can be seen
before cancellation:
\begin{align}
 \frac{\dd\cA_1}{V_\phi}
 &=\underbrace{\frac12}_{\rm stack}
 \underbrace{\frac{\dd^2\tau}{\tau_2}}_{\rm modulus/CKV}
 \underbrace{|\eta|^4}_{bc}
 \underbrace{\frac{R_{\rm phys}}
 {\sqrt{\A\tau_2}|\eta|^2}\Theta_1}_{X}
 \underbrace{\frac1{2\pi\sqrt{\A\tau_2}|\eta|^2}}_{L},
 \label{eq:genus1-factor-product}\\
 \Theta_1
 &:=\sum_{m,n\in\mathbb Z}
 \exp\!\left[-\frac{\pi R_{\rm phys}^2}{\A\tau_2}
 |m\tau-n|^2\right].
\end{align}
Thus $|\eta|^{4-2-2}=1$, the two factors $\tau_2^{-1/2}$ give
$\tau_2^{-1}$, and the moduli form supplies the second $\tau_2^{-1}$.  The
answer is
\begin{equation}
 \boxed{
 \frac1{V_\phi}\frac{\dd\cA_1}{\dd^2\tau}
 =\frac{R_{\rm phys}}{4\pi\A\tau_2^2}
 \sum_{m,n\in\mathbb Z}
 \exp\!\left[-\frac{\pi R_{\rm phys}^2}{\A\tau_2}
 |m\tau-n|^2\right].}
 \label{eq:torus-integrand}
\end{equation}
This is the genus-one integrand used by the direct $\tau$-plane computation.
There is no sphere-topology correction because $\Lambda_1=1$.

\subsection{Unfolding check and the universal logarithm}

The zero orbit uses
$\int_{\cF_1}\dd^2\tau/\tau_2^2=\pi/3$ and contributes
$R_{\rm phys}/(12\A)$.  Every nonzero pair can be written as a primitive pair
times $k\geq1$; modular images unfold $\cF_1$ to the strip.  Hence
\begin{align}
 \left.\frac{\cA_1}{V_\phi}\right|_{\rm nonzero}
 &=\frac{R_{\rm phys}}{4\pi\A}\,2
 \sum_{k=1}^\infty
 \int_{-1/2}^{1/2}\dd\tau_1
 \int_0^\infty\frac{\dd\tau_2}{\tau_2^2}
 \ee^{-\pi R_{\rm phys}^2k^2/(\A\tau_2)} \notag\\
 &=\frac{R_{\rm phys}}{4\pi\A}\,2
 \sum_{k=1}^\infty\frac{\A}{\pi R_{\rm phys}^2k^2}
 =\frac1{12R_{\rm phys}}.
\end{align}
Therefore
\begin{equation}
 \boxed{
 \frac{\cA_1}{V_\phi}
 =\frac1{12}\left(\frac{R_{\rm phys}}{\A}
 +\frac1{R_{\rm phys}}\right).}
 \label{eq:torus-result}
\end{equation}
The Liouville potential depends on its zero mode through
$\mu_L\exp(2bX^1/\sqrt\A)$.  Shifting the zero mode to absorb $\mu_L$ gives
\begin{equation}
 V_\phi^{\rm ren}
 =-\frac{\sqrt\A}{2b}\log\mu_L+C_{\rm cutoff}.
 \label{eq:liouville-volume}
\end{equation}
At $b=1$ and $\A=1$, where $r=R_{\rm phys}$,
\begin{equation}
 \cA_1\big|_{\log\mu_L}
 =-\frac1{24}\left(r+\frac1r\right)\log\mu_L.
 \label{eq:torus-log}
\end{equation}

\section{Genus two: derivation of every ingredient}

\subsection{Period-matrix coordinates and the ghost top form}

Let
\begin{equation}
 \Omega=X+\ii Y=\begin{pmatrix}\Omega_{11}&\Omega_{12}\\
 \Omega_{12}&\Omega_{22}\end{pmatrix},
 \qquad Y\succ0,
\end{equation}
and order the three holomorphic differentials on moduli space as
\begin{equation}
 \dd^3\Omega:=\dd\Omega_{11}\wedge\dd\Omega_{22}
 \wedge\dd\Omega_{12}.
 \label{eq:d3omega-order}
\end{equation}
There are three $b$-ghost zero modes because
$\dim_{\mathbb C}\cM_2=3$.  In the basis
$\omega_1^2,\omega_2^2,\omega_1\omega_2$ of holomorphic quadratic
differentials, contraction with the three Beltrami differentials converts the
ghost zero-mode determinant into $\dd^3\Omega$.  This is the content of Xi's
equations (4.100)--(4.103); it explains why no separate insertion-point
Jacobian remains in the vacuum measure.

For one complex coordinate $z=x+\ii y$,
$\dd z\wedge\dd\bar z=-2\ii\dd x\wedge\dd y$.  Reordering the three barred
forms next to their unbarred partners takes three transpositions.  Therefore
\begin{equation}
 \dd^3\Omega\wedge\dd^3\bar\Omega
 =(-1)^3(-2\ii)^3\,\dd^3X\dd^3Y
 =-8\ii\,\dd^3X\dd^3Y.
 \label{eq:six-form-jacobian-full}
\end{equation}
The factor $8$ is solely a complex-to-real form Jacobian.

\subsection{Critical seed and the 26 loop-momentum Gaussians}

Before doing the harmonic-momentum integrals, Xi's equation (4.105) is
\begin{equation}
 f(\Sigma)Z_X^{26}(\Sigma)
 =V_X\frac{g_s^2\A}{8\pi}
 \int\prod_{I=1}^{2}\frac{\dd^{26}k_I}{(2\pi)^{26}}
 \ee^{-\pi\A k_I\cdot k_JY_{IJ}}
 |\Phi(\Omega)|^2.
 \label{eq:xi-4105-full}
\end{equation}
For one target component, diagonalizing the positive matrix $Y$ gives
\begin{equation}
 \int\frac{\dd^2k}{(2\pi)^2}\ee^{-\pi\A k^TYk}
 =\frac1{(2\pi)^2}\frac1{\sqrt{\det(\A Y)}}
 =\frac1{4\pi^2\A\sqrt{\det Y}}.
 \label{eq:two-loop-gaussian-full}
\end{equation}
There are 26 components, hence
$(4\pi^2\A)^{-26}(\det Y)^{-13}$.  Rotating the two timelike loop momenta
to Euclidean contours gives $\ii^2=-1$.  The complex six-form is therefore
\begin{equation}
 \Omega_6^{\rm crit}
 =-V_X\frac{g_s^2\A}{8\pi}
 \frac{|\Phi(\Omega)\dd^3\Omega|^2}
 {(4\pi^2\A)^{26}(\det Y)^{13}}.
 \label{eq:critical-six-form-full}
\end{equation}

\begin{warning}
The factor $g_s^2\A/(8\pi)$ displayed in Xi's equation (4.105) must remain
after the Gaussian.  It is absent from the printed equation (4.106), but the
preceding equation and the gauge-fixing recurrence require it.
\end{warning}

\subsection{Oscillators, the Igusa form, and the factor
\texorpdfstring{$2^{24}$}{2 to the 24}}

The combined chiral oscillator function of the ghosts and 26 scalars has
modular weight $-10$ and unit plumbing residue.  It is uniquely
\begin{equation}
 \Phi(\Omega)=\frac{2\pi\ii}{\chi_{10}^{\rm note}(\Omega)}.
 \label{eq:Phi-note}
\end{equation}
To compare this with the implementation, define the genus-two theta constant
\begin{equation}
 \theta\!\begin{bmatrix}a\\b\end{bmatrix}(0|\Omega)
 =\sum_{\ell\in\mathbb Z^2}
 \exp\!\left[
 \pi\ii(\ell+a)^T\Omega(\ell+a)+2\pi\ii(\ell+a)^Tb\right]
\end{equation}
and the raw product over the ten even characteristics,
\begin{equation}
 \Psi_{10}(\Omega)=\prod_{\delta\ {\rm even}}
 \theta[\delta](0|\Omega)^2,
 \qquad
 \chi_{10}^{\rm note}=2^{-12}\Psi_{10}.
 \label{eq:chi-conventions-full}
\end{equation}
Let $\epsilon=\Omega_{12}$.  With the literal plumbing coordinate,
\begin{equation}
 q=2\pi\ii\epsilon+O(\epsilon^3),
 \qquad
 \Psi_{10}=2^{12}q^2\eta(\Omega_{11})^{24}
 \eta(\Omega_{22})^{24}+O(q^4).
 \label{eq:Psi-separating}
\end{equation}
Consequently $\chi_{10}^{\rm note}\sim
q^2\eta_1^{24}\eta_2^{24}$ and
$\dd\epsilon=\dd q/(2\pi\ii)$.  Substitution into
\eqref{eq:Phi-note} gives the unit chiral residue
\begin{equation}
 \Phi\,\dd\Omega_{11}\dd\Omega_{22}\dd\epsilon
 \longrightarrow
 \frac{\dd q}{q^2}
 \frac{\dd\Omega_{11}}{\eta(\Omega_{11})^{24}}
 \frac{\dd\Omega_{22}}{\eta(\Omega_{22})^{24}}.
 \label{eq:unit-mumford-residue}
\end{equation}
If one instead evaluates $(2\pi\ii)/\Psi_{10}$, its chiral residue is
$2^{-12}$.  Its absolute square must therefore be multiplied by
\begin{equation}
 \boxed{(2^{12})^2=2^{24}.}
 \label{eq:raw-product-conversion}
\end{equation}
This factor belongs to the local Mumford density; it is unrelated to either
the sphere constant or the moduli-stack quotient.

\subsection{The compact scalar at genus two}

Normalize the holomorphic one-forms by
$\oint_{A_I}\omega_J=\delta_{IJ}$ and
$\oint_{B_I}\omega_J=\Omega_{IJ}$.  A harmonic map to the circle is labelled
by $m,n\in\mathbb Z^2$.  Using
$\ii\int_\Sigma\omega_I\wedge\bar\omega_J=2Y_{IJ}$ in the scalar action gives
\begin{equation}
 S_{m,n}^{(2)}
 =\pi r^2(m+\bar\Omega n)^TY^{-1}(m+\Omega n),
 \qquad r=\frac{R_{\rm phys}}{\sqrt\A}.
 \label{eq:genus2-classical-action}
\end{equation}
Thus the classical lattice sum is
\begin{equation}
 \boxed{
 \Theta_r^{(2)}(\Omega)
 =\sum_{m,n\in\mathbb Z^2}
 \exp\!\left[-\pi r^2
 (m+\bar\Omega n)^TY^{-1}(m+\Omega n)\right].}
 \label{eq:theta-genus2-full}
\end{equation}
The connected compact zero mode is the target length, not merely the radius:
\begin{equation}
 \int_0^{2\pi R_{\rm phys}}\dd x_0=2\pi R_{\rm phys}
 =2\pi\sqrt\A\,r.
 \label{eq:genus2-compact-zero-mode}
\end{equation}
If $Z_X^\Xi$ is one noncompact scalar partition per unit target length in
Xi's convention, the compact scalar is therefore
\begin{equation}
 Z_{X,r}^\Xi=(2\pi\sqrt\A\,r)Z_X^\Xi\Theta_r^{(2)}.
 \label{eq:compact-from-noncompact}
\end{equation}

\subsection{Liouville sewing and the Xi-normalized same-frame quotient}

In a three-tube pants decomposition $a$, the normalized Liouville vacuum
partition is defined entirely by the input two- and three-point data:
\begin{equation}
 Z_L^{(a)}
 =\int_0^\infty\prod_{e=1}^3\frac{\dd P_e}{\pi}
 C(P_1,P_2,P_3)^2
 \left|\cF_{25}^{(a)}(\boldsymbol P;\boldsymbol q)\right|^2.
 \label{eq:liouville-genus2-sewing}
\end{equation}
The Virasoro block includes all descendant Gram matrices and the powers of
$q_e$ dictated by the weights.  Equation \eqref{eq:liouville-genus2-sewing}
shows explicitly that there is no extra Liouville normalization constant:
each internal edge carries exactly the completeness measure $\dd P/\pi$ and
each pair of pants carries one DOZZ coefficient.

\paragraph{Vacuum shifts and the meaning of $A_L$.}
Let $\cB_X^{(a)}$ and $\cB_L^{(a)}$ denote the unshifted scalar and
Liouville blocks in the same three-tube plumbing channel, and define
\begin{equation}
 \cV_a:=\prod_{e=1}^3|q_e|^{-1/12}.
\end{equation}
The Xi-normalized partitions are
\begin{equation}
 Z_X^{\Xi,(a)}=A_X\cV_a\cB_X^{(a)},
 \qquad
 Z_L^{\Xi,(a)}=A_L^{\rm sewn}\cV_a^{25}\cB_L^{(a)},
 \qquad
 A_X=\frac1{4\pi^2\A},
 \quad A_L^{\rm sewn}=1.
 \label{eq:AL-sewn-ledger}
\end{equation}
Here $A_X$ is Xi's two-handle momentum-measure normalization, whereas
$A_L^{\rm sewn}=1$ follows from using literally the primary, metric, inverse
metric, DOZZ coefficient, and intrinsic momentum $P$ of Appendix~I.  Hence
\begin{equation}
 \frac{Z_L^{\Xi,(a)}}{(Z_X^{\Xi,(a)})^{25}}
 =(4\pi^2\A)^{25}
 \frac{\cB_L^{(a)}}{(\cB_X^{(a)})^{25}}.
 \label{eq:q-shift-cancellation-quotient}
\end{equation}
The factor $\cV_a^{25}$ cancels exactly.  Thus the
$|q_e|^{-25/12}$ shifts in the Liouville block produce no further constant;
the surviving $(4\pi^2\A)^{25}$ is the scalar normalization and must be
combined with the critical $26$-scalar seed and compact zero mode.

\paragraph{What genus one does and does not determine.}
Xi's equation (4.92) writes the unpunctured torus trace per ordinary
Liouville length as
\begin{equation}
 \frac{Z_L^{\Xi,(1)}}{V_\phi}
 =\int_0^\infty\frac{\dd P_{\rm phys}}{\pi}
 \frac{\ee^{-\pi\A\tau_2P_{\rm phys}^2}}{|\eta(\tau)|^2}
 =\frac1{2\pi\sqrt{\A\tau_2}|\eta(\tau)|^2}.
 \label{eq:Xi-Liouville-torus-volume-density}
\end{equation}
The intrinsic CFT momentum of Appendix~I instead gives
\begin{equation}
 \cZ_L^{\rm CFT,(1)}
 =\int_0^\infty\frac{\dd P_{\rm CFT}}{\pi}
 \frac{\ee^{-4\pi\tau_2P_{\rm CFT}^2}}{|\eta(\tau)|^2}
 =\frac1{4\pi\sqrt{\tau_2}|\eta(\tau)|^2}.
 \label{eq:CFT-Liouville-torus-density}
\end{equation}
Matching only the Gaussian characters gives
\begin{equation}
 P_{\rm phys}=\frac{2P_{\rm CFT}}{\sqrt\A},
 \qquad
 \frac{Z_L^{\Xi,(1)}}{V_\phi}
 =\underbrace{\frac2{\sqrt\A}}_{J_\phi^{(1)}}
 \cZ_L^{\rm CFT,(1)}.
 \label{eq:Jphi-genus-one}
\end{equation}
The factor $J_\phi^{(1)}$ converts the dimensionful asymptotic wave number
and the regulator defining the divergent ordinary length $V_\phi$.  It is
not $A_L^{\rm sewn}$ and must not be inserted into the genus-two DOZZ
integral.  The symbol $P$ is overloaded in the note: the DOZZ arguments in
(4.119), the weight $h=1+P^2$, and every momentum in
\eqref{eq:liouville-genus2-sewing} remain the dimensionless intrinsic
$P_{\rm CFT}$.  Replacing them by $2P_{\rm CFT}/\sqrt\A$ would change the
arguments and weights of the non-homogeneous DOZZ function.  Therefore
\begin{equation}
 \boxed{J_\phi^{(1)}=\frac2{\sqrt\A},
 \qquad A_L^{\rm sewn}=1.}
 \label{eq:AL-final-verdict}
\end{equation}

For the analytic derivation, place $Z_L^{(a)}$ together with Xi's noncompact
scalar $Z_{X,k}^{\Xi,(a)}$ in the same conformal frame.  Their Weyl
transformations have
the form
\begin{equation}
 Z_{X,k}^{\Xi,(a)}=W_a Z_{X,k}^{\Xi,B},
 \qquad
 Z_L^{(a)}=W_a^{25}Z_L^{B},
\end{equation}
so the combination
\begin{equation}
 \frac{Z_L^{(a)}}{(Z_{X,k}^{\Xi,(a)})^{25}}
 =\frac{Z_L^{B}}{(Z_{X,k}^{\Xi,B})^{25}}
 \label{eq:same-frame-quotient}
\end{equation}
is frame independent.  This is why a separately guessed Weyl determinant
must not be inserted into the integrand.  The repository convention will be
introduced only after the Xi-normalized density has been derived.

\subsection{Xi-normalized momentum measure and matter replacement}

Xi's note writes the physical momentum as $k$ and fixes the free-boson states
by
\begin{equation}
 \boxed{
 \langle k|k'\rangle=2\pi\delta(k-k'),
 \qquad
 \mathbf 1=\int_{\mathbb R}\frac{\dd k}{2\pi}|k\rangle\langle k|.}
 \label{eq:xi-k-state-normalization}
\end{equation}
This is the convention stated explicitly in Appendix P around (P.20).
Appendix E gives the general form $C_X2\pi\delta(k-k')$; Xi's convention here
sets $C_X=1$.  All formulas through \eqref{eq:I2Xi-full} use this convention.

The $k_I$ are the two handle momenta, not the connected constant target-space
mode $x_0$.  The latter is divided out for a noncompact scalar and becomes the
target length $2\pi R_{\rm phys}=2\pi\sqrt\A\,r$ for a compact scalar.  Write
\begin{equation}
 Y=\begin{pmatrix}a&b\\b&c\end{pmatrix},
 \qquad D:=\det Y=ac-b^2>0.
\end{equation}
Xi's handle-momentum integral for one scalar per ordinary target-space volume
is therefore
\begin{align}
 I_{0,k}^{\Xi}(Y)
 &:=\int_{\mathbb R^2}\frac{\dd k_1}{2\pi}
 \frac{\dd k_2}{2\pi}
 \exp\!\left[-\pi\A(ak_1^2+2bk_1k_2+ck_2^2)\right] \notag\\
 &=\frac1{(2\pi)^2}
 \int\dd k_1\dd k_2\,
 \exp\!\left[-\pi\A a\left(k_1+\frac ba k_2\right)^2
 -\pi\A\frac Da k_2^2\right] \notag\\
 &=\boxed{\frac1{4\pi^2\A\sqrt{\det Y}}}.
 \label{eq:k-zero-mode-integral}
\end{align}

For bookkeeping inside the same Xi normalization, change variables only:
$p_I=\sqrt\A\,k_I$.  Define
\begin{equation}
 I_{0,p}^{(2\pi)}(Y)
 :=\int_{\mathbb R^2}\frac{\dd p_1}{2\pi}
 \frac{\dd p_2}{2\pi}\,\ee^{-\pi p^TYp}
 =\frac1{4\pi^2\sqrt{\det Y}}.
 \label{eq:xi-p-zero-mode-integral}
\end{equation}
The genuine genus-two change of variables is then
\begin{equation}
 \prod_{I=1}^2\frac{\dd k_I}{2\pi}
 =\frac1\A\prod_{I=1}^2\frac{\dd p_I}{2\pi},
 \qquad
 \boxed{Z_{X,k}^\Xi=\frac1\A Z_{X,p}^{(2\pi)}}.
 \label{eq:scalar-conversion-full}
\end{equation}
No extra $2\pi$ occurs in this $k\to p$ map: Xi's factor $1/(2\pi)$ is
present on both sides.  If actual dimensionless-$p$ states are introduced,
they must be rescaled as
$|p\rangle=\A^{-1/4}|k=p/\sqrt\A\rangle$ in order to retain
$\langle p|p'\rangle=2\pi\delta(p-p')$.

Now perform the matter replacement without referring to the repository:
\begin{equation}
 Z_X^{26}\longrightarrow Z_{X,r}^\Xi Z_L,
 \qquad
 Z_{X,r}^\Xi=(2\pi\sqrt\A\,r)Z_{X,k}^\Xi\Theta_r^{(2)},
 \label{eq:c1-replacement-full}
\end{equation}
and hence
\begin{equation}
 \frac{Z_{X,r}^\Xi Z_L}{(Z_{X,k}^\Xi)^{26}}
 =(2\pi\sqrt\A\,r)\Theta_r^{(2)}
 \frac{Z_L}{(Z_{X,k}^\Xi)^{25}}.
 \label{eq:c1-replacement-quotient}
\end{equation}
Using the raw theta product and temporarily omitting the external topology
coefficient, the Xi-normalized local density is
\begin{align}
 \cD_{2,\rm raw}^{c=1,\Xi}
 &=\underbrace{\frac{|2\pi\ii|^2}
 {(4\pi^2\A)^{26}(\det Y)^{13}|\Psi_{10}|^2}}
 _{\text{critical density in Xi's }\dd k/(2\pi)\text{ convention}}
 \underbrace{(2\pi\sqrt\A\,r)\Theta_r^{(2)}}
 _{\text{compact zero mode and windings}}
 \underbrace{\frac{Z_L}{(Z_{X,k}^\Xi)^{25}}}
 _{\text{Xi-normalized same-frame quotient}} \notag\\
 &=\boxed{
 \frac{2\pi r}{\sqrt\A}(4\pi^2)^{-25}
 \frac{\Theta_r^{(2)}}{(\det Y)^{13}|\Psi_{10}|^2}
 \frac{Z_L}{(Z_{X,p}^{(2\pi)})^{25}}}.
 \label{eq:raw-local-density}
\end{align}
The second line merely uses the Xi-preserving change of variables
\eqref{eq:scalar-conversion-full}:
\begin{equation}
 |2\pi\ii|^2(4\pi^2\A)^{-26}
 \A^{25}(2\pi\sqrt\A\,r)
 =\frac{2\pi r}{\sqrt\A}(4\pi^2)^{-25}.
 \label{eq:correlated-scalar-conversion-full}
\end{equation}
Equivalently, the 26 critical scalars and the inverse 25-scalar quotient leave
one genuine Xi $k\to p$ Jacobian,
\begin{equation}
 (\A^{-1})^{26}(\A^{-1})^{-25}=\A^{-1}.
 \label{eq:26-minus-25}
\end{equation}

Applying the unit-residue conversion \eqref{eq:raw-product-conversion} now
defines the integrand before any code convention is mentioned:
\begin{align}
 \cI_2^\Xi(\Omega;r,\A)
 &:=\boxed{2^{24}\frac{|2\pi\ii|^2}{(4\pi^2\A)^{26}}
 \frac{(2\pi\sqrt\A\,r)\Theta_r^{(2)}(\Omega)}
 {(\det Y)^{13}|\Psi_{10}(\Omega)|^2}
 \frac{Z_L(\Omega)}{(Z_{X,k}^\Xi(\Omega))^{25}}} \notag\\
 &=2^{24}\frac{2\pi r}{\sqrt\A}(4\pi^2)^{-25}
 \frac{\Theta_r^{(2)}(\Omega)}
 {(\det Y)^{13}|\Psi_{10}(\Omega)|^2}
 \frac{Z_L(\Omega)}{(Z_{X,p}^{(2\pi)}(\Omega))^{25}}.
 \label{eq:I2Xi-full}
\end{align}

\paragraph{Translation to the repository evaluator.}
Only now introduce the repository's independent bare-measure object:
\begin{align}
 I_{0,p}^{\rm code}(Y)
 &:=\int_{\mathbb R^2}\dd p_1\dd p_2\,\ee^{-\pi p^TYp}
 =\frac1{\sqrt{\det Y}}
 =(2\pi)^2 I_{0,p}^{(2\pi)}(Y),
 \label{eq:p-zero-mode-integral}\\
 Z_{X,p}^{\rm code}&=(2\pi)^2Z_{X,p}^{(2\pi)}.
 \label{eq:code-vs-xip-scalar}
\end{align}
The oscillator determinant is common to both objects.  Combining this
repository-only conversion with the Xi $k\to p$ Jacobian gives
\begin{equation}
 \boxed{
 Z_{X,k}^\Xi
 =\frac1\A Z_{X,p}^{(2\pi)}
 =\frac1{4\pi^2\A}Z_{X,p}^{\rm code}.}
 \label{eq:code-to-xi-scalar}
\end{equation}
For clarity, the two factors are
\begin{equation}
 C_\A:=\A^{-1},
 \qquad C_{2\pi}:=(2\pi)^{-2},
 \qquad C_{\mathrm{code}\to k}:=C_\A C_{2\pi}
 =\frac1{4\pi^2\A}.
 \label{eq:two-scalar-conversions}
\end{equation}
Here $C_\A$ belongs to the Xi derivation; $C_{2\pi}$ belongs only to the final
translation of the repository object.  Substitution into
\eqref{eq:I2Xi-full} gives the numerically convenient but logically secondary
form
\begin{equation}
 \boxed{
 \cI_2^\Xi
 =2^{24}\frac{2\pi r}{\sqrt\A}
 \frac{\Theta_r^{(2)}}{(\det Y)^{13}|\Psi_{10}|^2}
 \frac{Z_L}{(Z_{X,p}^{\rm code})^{25}}.}
 \label{eq:I2Xi-xip-form}
\end{equation}

\subsection{Gauge phase, sphere topology, and the final real coefficient}

The critical six-form contains four independent factors: the phase
$N_{2,0}=-\ii$, the Wick-rotation sign $-1$, the form Jacobian $-8\ii$, and
the coefficient $g_s^2\A/(8\pi)$.  Their product is
\begin{equation}
 (-\ii)(-1)(-8\ii)\frac{g_s^2\A}{8\pi}
 =\frac{g_s^2\A}{\pi}.
 \label{eq:critical-positive-full}
\end{equation}
After stripping $g_s^2$, the critical topology coefficient is $\A/\pi$.
Replacing its sphere metric by the reduced physical $c=1$ metric multiplies
it by \eqref{eq:Lambda12}, hence
\begin{equation}
 \underbrace{\frac{\A}{\pi}}_{\rm critical\ real\ form}
 \underbrace{\frac2\A}_{\rm sphere\ sewing}
 =\boxed{\frac2\pi}.
 \label{eq:final-external-coefficient-full}
\end{equation}
For comparison with the historical dimensionless-plumbing implementation,
let $\cK_2^{\rm old}$ include its original critical coefficient $\A/\pi$,
its dimensionless compact prefactor $r$, and the raw same-frame quotient
$\cB_L/(\cB_X)^{25}$.  The correlated Xi scalar conversion and the separate
topology replacement are
\begin{equation}
 \cN_{\rm scalar}=\frac1{2\pi\sqrt\A},
 \qquad
 \cN_{\rm top}=\frac2\A,
 \qquad
 \boxed{\cN_{\rm final}
 =\cN_{\rm scalar}\cN_{\rm top}
 =\frac1{\pi\A^{3/2}}.}
 \label{eq:final-normalization-relative-old}
\end{equation}
Thus $\cK_2^{c=1}=\cN_{\rm final}\cK_2^{\rm old}$ in this explicitly
defined comparison, and $\cN_{\rm final}=1/\pi$ at $\A=1$.  This relative
multiplier must not be confused with the local coefficient $2/\pi$ below:
the latter multiplies the already Xi-normalized object $\cI_2^\Xi$, whereas
$\cN_{\rm final}$ compares directly to the pre-conversion historical
dimensionless-plumbing expression.
The final positive-real local kernel is therefore
\begin{equation}
 \boxed{
 K_2^{c=1}(\Omega;r,\A)
 =\frac2\pi\cI_2^\Xi(\Omega;r,\A).}
 \label{eq:final-kernel-full}
\end{equation}
Equivalently, after the separate repository translation
\eqref{eq:code-to-xi-scalar}, the evaluator may display every factor as
\begin{equation}
 \boxed{
 K_2^{c=1}
 =\underbrace{\frac2\pi}_{\rm topology/form}
 \underbrace{2^{24}}_{\rm raw\ \Psi_{10}}
 \underbrace{\frac{2\pi r}{\sqrt\A}}_{\rm scalar\ measures}
 \underbrace{\Theta_r^{(2)}}_{\rm compact\ lattice}
 \underbrace{(\det Y)^{-13}}_{26\ \rm Gaussians}
 \underbrace{|\Psi_{10}|^{-2}}_{\rm oscillators}
 \underbrace{Z_L (Z_{X,p}^{\rm code})^{-25}}_{\rm CFT\ replacement}.}
 \label{eq:factorized-final-kernel}
\end{equation}

When a coarse $PSp(4,\mathbb Z)$ fundamental domain is used, a generic
genus-two curve has the hyperelliptic involution.  Dividing by this stabilizer
gives a separate global weight $1/2$:
\begin{equation}
 \boxed{
 \frac{\cA_2(r)}{g_s^2}
 =\frac12\int_{\cF_2^{\rm coarse}}
 \dd^3X\dd^3Y\,K_2^{c=1}
 =\frac1\pi\int_{\cF_2^{\rm coarse}}
 \dd^3X\dd^3Y\,\cI_2^\Xi.}
 \label{eq:final-amplitude-full}
\end{equation}
The local $2/\pi$ and the global $1/2$ must each be applied exactly once.

\subsection{Numerical reference at the self-dual radius}

For reference, set $\A=1$ and $R_{\rm phys}=r=1$.  The latest completed
physical-mixture RQMC run used eight independent scrambled-Sobol replicates,
four proposal components, and $256$ proposals per component; $6846$ in-domain
nodes required CFT evaluation.  The archived table used the legacy
dimensionless-scalar density $\cI_2^{\rm old}$ and the label
\begin{equation}
 K_{2,\rm legacy}^{\rm note}=\frac1\pi\cI_2^{\rm old},
 \qquad
 \cI_2^\Xi=\frac1{2\pi}\cI_2^{\rm old}
             =\frac12K_{2,\rm legacy}^{\rm note}
\end{equation}
and found
\begin{equation}
 \int_{\cF_2^{\rm coarse}}\dd^3X\dd^3Y\,
 K_{2,\rm legacy}^{\rm note}(\Omega;1)
 =\left(3.4977126795939466
 \mathbin{\pm}0.08214314738188904\right)\times10^{-4}.
 \label{eq:R1-saved-Knote-integral}
\end{equation}
After the exact scalar-measure conversion, the direct numerical reference for
the Xi-normalized integrand
defined in \eqref{eq:I2Xi-full} is
\begin{equation}
 \boxed{
 \int_{\cF_2^{\rm coarse}}\dd^3X\dd^3Y\,
 \cI_2^\Xi(\Omega;1)
 =\left(1.7488563397969733
 \mathbin{\pm}0.04107157369094452\right)\times10^{-4}.}
 \label{eq:R1-I2Xi-integral}
\end{equation}
The uncertainty is the standard error across the eight complete replicates.
Equation \eqref{eq:R1-I2Xi-integral} is a coarse-domain integral and does not
include the generic stack weight $1/2$.  Combining it with the final local
coefficient $2/\pi$ and that stack weight gives
\begin{equation}
 \left.\frac{\cA_2}{g_s^2}\right|_{r=1}
 =\frac1\pi\int_{\cF_2^{\rm coarse}}\dd^3X\dd^3Y\,\cI_2^\Xi
 =\left(0.5566782624725754
 \mathbin{\pm}0.013073487946953723\right)\times10^{-4}.
 \label{eq:R1-final-amplitude-reference}
\end{equation}
This is the release field labelled
\texttt{connected\_logZ\_genus2\_over\_gs\_squared}.  The corresponding
coarse integral of $K_2^{c=1}$ is twice this value.  No matrix-model
normalization is used in any of these numbers.  The machine-readable source is
\path{output/data/genus2_c1_free_energy_direct_39/summary.json}.

\subsection{One-to-one map to the numerical evaluator}

At $\A=1$ the public code radius is $r$.  The returned quantities correspond
to the analytic factors as follows:
\begin{center}
\begin{tabular}{@{}p{0.37\textwidth}p{0.52\textwidth}@{}}
\toprule
Implementation quantity & Analytic factor \\
\midrule
\texttt{compact\_target\_zero\_mode} & $2\pi r$ \\
\texttt{compact\_winding\_sum} & $\Theta_r^{(2)}$ \\
\texttt{det\_im\_omega} & $\det Y$ \\
\texttt{chi10} with \texttt{product} & $\Psi_{10}$ \\
\texttt{liouville\_partition} & $Z_L$ from \eqref{eq:liouville-genus2-sewing} \\
\texttt{noncompact\_scalar\_partition} & $Z_{X,p}^{\rm code}$ with bare $\dd^2p$; converted to Xi's physical-$k$ object by \eqref{eq:code-to-xi-scalar} \\
\texttt{factorization\_normalization} & $2^{24}$ \\
historical critical multiplier & $\A/\pi$ before the $c=1$ sphere replacement \\
\texttt{c1\_genus2\_topology\_correction} & $2/\A$ \\
\texttt{string\_note\_kernel\_multiplier} & final sphere-normalized product $2/\pi$ \\
Monte Carlo stack weight & $1/2$ \\
\bottomrule
\end{tabular}
\end{center}
The v3 production evaluator applies the scalar-measure conversion and the
sphere-normalized multiplier exactly once.  When archived values are read,
their convention identifier selects an exact migration before integration;
the stack weight remains a separate Monte Carlo factor.


\section{Factor ledger}

\subsection*{Genus one}
\begin{center}
\small
\begin{tabular}{@{}p{0.29\textwidth}p{0.22\textwidth}p{0.40\textwidth}@{}}
\toprule
Factor & Value & Origin \\
\midrule
$X$ zero mode & $R_{\rm phys}/\sqrt\A$ & \eqref{eq:compact-zero-mode}; also Poisson resummation \\
$X$ classical action & $\pi R_{\rm phys}^2|m\tau-n|^2/(\A\tau_2)$ & action and torus metric \\
$X$ oscillators & $\tau_2^{-1/2}|\eta|^{-2}$ & scalar determinant \\
Liouville trace per $V_\phi$ & $(2\pi\sqrt{\A\tau_2}|\eta|^2)^{-1}$ & $\dd P/\pi$ half-line Gaussian \\
ghost oscillators & $|\eta|^4$ & left-right $bc$ determinant \\
modulus/CKV & $\dd^2\tau/\tau_2$ & \eqref{eq:torus-real-measure} \\
torus stack weight & $1/2$ & generic $z\mapsto-z$ stabilizer \\
$N_{1,0}$ and topology & $1$ & plumbing unitarity and $\Lambda_1=1$ \\
\bottomrule
\end{tabular}
\end{center}

\subsection*{Genus two}
\begin{center}
\small
\begin{tabular}{@{}p{0.29\textwidth}p{0.22\textwidth}p{0.40\textwidth}@{}}
\toprule
Factor & Value & Origin \\
\midrule
$N_{1,0}$ & $1$ & Plumbing unitarity, \eqref{eq:critical-solution} \\
$N_{2,0}$ & $-\ii$ & Plumbing unitarity, \eqref{eq:critical-solution} \\
$K_{S^2}^{\rm crit}$ & $8\pi/\A$ & BRST state metric and pole residue \\
$\wt K_{S^2}$ & $2/\sqrt{\A}$ & Physical $c=1$ state unitarity \\
$\wh K_{S^2}^{c=1}$ & $4\pi$ & Fourier factor and \eqref{eq:delta-jacobian} \\
$\Lambda_2$ & $2/\A$ & Sphere topology sewing \\
one-scalar Gaussian & $(4\pi^2\A\sqrt{\det Y})^{-1}$ & \eqref{eq:two-loop-gaussian-full} \\
26-scalar Gaussian & $(4\pi^2\A)^{-26}(\det Y)^{-13}$ & product of 26 Gaussians \\
Mumford numerator & $|2\pi\ii|^2=4\pi^2$ & unit $q$-residue in Xi's convention \\
raw $\Psi_{10}$ conversion & $2^{24}$ & \eqref{eq:raw-product-conversion} \\
compact zero mode & $2\pi\sqrt\A\,r$ & target-circle length \\
compact lattice & $\Theta_r^{(2)}$ & harmonic-map action \\
Xi state completeness & $\dd k/(2\pi)$ & \eqref{eq:xi-k-state-normalization} \\
Xi $k\to p$ Jacobian & $\A^{-1}$ & \eqref{eq:scalar-conversion-full} \\
bare code $p\to(2\pi)$-measure $p$ & $(2\pi)^{-2}$ & \eqref{eq:code-vs-xip-scalar} \\
combined code $p\to$ Xi $k$ & $(4\pi^2\A)^{-1}$ & \eqref{eq:code-to-xi-scalar} \\
code-form scalar product & $2\pi r/\sqrt\A$ & \eqref{eq:I2Xi-xip-form} \\
complex six-form Jacobian & $-8\ii$ & \eqref{eq:six-form-jacobian-full} \\
final local coefficient & $2/\pi$ & \eqref{eq:final-external-coefficient-full} \\
genus-two stack weight & $1/2$ & generic hyperelliptic involution \\
\bottomrule
\end{tabular}
\end{center}

\section{Matrix model only as an external test}

Nothing above uses the matrix model.  After the worldsheet normalization has
been fixed, the tree amplitude \eqref{eq:tree-amplitude} may be compared with
the collective-field result
\begin{equation}
 \cA_{1\to2}^{\rm MQM}
 =\frac{\ii}{\mu}\omega\omega_1\omega_2.
\end{equation}
This gives the dictionary
\begin{equation}
 \boxed{\mu^{-1}=4\pi g_s^{\Xi}.}
 \label{eq:mu-gs}
\end{equation}
If
\begin{equation}
 f_2(r)=\frac{7r^2+10+7r^{-2}}{5760r},
\end{equation}
then the matrix-model prediction in the same coupling convention is
\begin{equation}
 \frac{\cA_2^{\rm MM}(r)}{(g_s^\Xi)^2}
 =16\pi^2 f_2(r).
 \label{eq:mm-target}
\end{equation}
Equation \eqref{eq:mm-target} is a test of \eqref{eq:final-amplitude-full}, not an
input used to alter it.

\section{Items that still require an independent review}

The derivation reduces the absolute normalization question to the following
sharp checks.
\begin{enumerate}
\item Verify directly from Xi's BPZ and physical-state conventions that the
conversion from \eqref{eq:Ktilde} to \eqref{eq:Khat} contains exactly one
Fourier $2\pi$ and exactly one Jacobian $\sqrt{\A}$.
\item Verify that the scalar conversion \eqref{eq:correlated-scalar-conversion-full}
is included once and is not mistaken for the topology replacement
\eqref{eq:Lambda12} merely because both contain related zero-mode constants.
\item Verify each global $1/2$ against the precise moduli stack used in the
numerical integration.
\item Test the analytic separating formula \eqref{eq:Psi-separating}
numerically with the literal coordinate $uv=q$ and
$q\simeq2\pi\ii\Omega_{12}$.
\end{enumerate}

If these checks pass, no additional moduli-independent factor is allowed by
unitarity or by a change of local sewing coordinates.

\section{Implementation manifest for the genus-two integral}

The exhaustive, reviewer-facing source inventory is kept in
\path{plumbing/genus2_monte_carlo_code_manifest.md}.  The numerical path is
not a single script.  It has the following auditable stages:
\begin{equation}
 \begin{aligned}
 &\text{physical-measure RQMC nodes}
 \longrightarrow
 \text{certified }(\Omega,q,\operatorname{Sp}(4,\mathbb Z))\text{ chart}
 \\
 &\hspace{2cm}\longrightarrow (Z_L,Z_X,\Theta_r,\Psi_{10})
 \longrightarrow K_2^{c=1}
 \longrightarrow \text{complete-scramble assembly}.
 \end{aligned}
\end{equation}
The source closure of the pointwise evaluator contains 31 local Python
modules.  The manifest lists every one, together with the period-table
builders, scheduler wrappers, radius reweighting and release code, and the
independent convention checks.  It also distinguishes code that changes the
value from optional plots, recovery jobs, and historical samplers.

For the preferred physical-period-measure design the implemented estimator is
\begin{equation}
 \left(\frac{\cA_2}{g_s^2}\right)_s
 =\frac12\frac1{4M}
 \sum_{a=1}^{4M}
 \mathbf 1_{\cF_2}(\Omega_{s,a})
 \frac{J_Y(t_{1,s,a},t_{3,s,a})}
      {p_{\rm mix}(t_{1,s,a},t_{3,s,a})}
 K_2^{c=1}(\Omega_{s,a};r).
 \label{eq:implemented-physical-mixture-estimator}
\end{equation}
Thus no $\det(Y)^3$ multiplies the local kernel in the current production
design.  Out-of-domain proposals contribute exact zero; every in-domain node
must have a certified CFT result; and the uncertainty is the standard error
across independent complete scrambles.  The exact control observable
$K_{\rm control}=\det(Y)^{-3}$ must reproduce
\begin{equation}
 \int_{\cF_2^{\rm coarse}}\frac{\dd^3X\dd^3Y}{\det(Y)^3}
 =\frac{\pi^3}{270}.
\end{equation}
These requirements are enforced by the design generator and strict assembler,
not added by post-processing.

\begin{thebibliography}{9}
\bibitem{XiNotes}
X.~Yin, \emph{String Theory Notes}, especially Sections 3.8, 4.6, 4.8--4.10
and Appendices E, I.5, and P; equations (4.32), (4.58)--(4.72),
(4.90)--(4.122), (I.39)--(I.42), and (P.20).

\bibitem{BRY}
B.~Balthazar, V.~A.~Rodriguez, and X.~Yin,
\emph{The $c=1$ String Theory S-Matrix Revisited},
\href{https://arxiv.org/abs/1705.07151}{arXiv:1705.07151}.

\bibitem{DHP}
E.~D'Hoker and D.~H.~Phong,
\emph{Two-Loop Superstrings VI: Non-Renormalization Theorems and the 4-Point
Function}, and related genus-two measure papers.
\end{thebibliography}

\end{document}
